% ==========================================
% SECCIÓN 2.3: COMPILACIÓN, EJECUCIÓN Y EJEMPLOS MPI
% RESPONSABLE: Jonatan
% DESCRIPCIÓN: Uso de `mpicc`, `mpirun`, explicación de `-np` y flujo básico.
% ==========================================

% * Explicación de mpicc.
% * Explicación de mpirun.
% * Significado de -np.
% * Flujo de compilación y ejecución.
% * Ejemplo sencillo en C usando MPI.
% * Explicación del resultado del programa.

\section{Compilación, ejecución y ejemplos en MPI}

En esta sección se abordarán los aspectos prácticos relacionados con la compilación y ejecución de programas escritos en C utilizando la biblioteca MPI. Se explicarán los comandos fundamentales como \texttt{mpicc} para la compilación y \texttt{mpirun} o \texttt{mpiexec} para la ejecución de aplicaciones paralelas. Además, se presentará un ejemplo ilustrativo de un programa de saludos distribuido, que permitirá comprender el flujo completo desde el código fuente hasta la salida en pantalla, destacando las funciones esenciales de MPI involucradas en el proceso.


\subsection{Explicación del comando \texttt{mpicc}}

El comando \texttt{mpicc} se utiliza para compilar programas escritos en C que emplean la biblioteca MPI \cite{pacheco_2011}. A diferencia de un compilador nativo, \texttt{mpicc} opera como un \textit{wrapper} que invoca al compilador subyacente del sistema, agregando automáticamente las rutas de inclusión de cabeceras como \texttt{mpi.h} y especificando las bibliotecas de enlace requeridas por el estándar \cite{pacheco_2011, mpi_forum_2025}. Esta abstracción permite al desarrollador evitar la configuración manual de dependencias complejas, facilitando la portabilidad del código fuente entre distintos clústeres y supercomputadoras \cite{grama_intro_parallel_2003}.

La sintaxis operativa de \texttt{mpicc} hereda la estructura convencional de los compiladores de C en sistemas tipo Unix, por lo que la invocación elemental consiste en indicar el comando, las banderas deseadas, los archivos fuente y el nombre del binario de salida mediante el modificador \texttt{-o} \cite{pacheco_2011}. Un caso típico de uso incluye las banderas de depuración y advertencias, como se muestra a continuación:

\begin{verbatim}
$ mpicc -g -Wall -o mpi_hello mpi_hello.c
\end{verbatim}

Como la mayoría de las opciones se transfieren directamente al compilador base subyacente (por ejemplo, \texttt{gcc}), el usuario puede combinar las banderas tradicionales de optimización o diagnóstico con la gestión automática de las bibliotecas de MPI \cite{pacheco_2011}. A fin de organizar las opciones más relevantes, en la Tabla~\ref{tab:mpicc_flags} se presentan los parámetros opcionales de compilación y diagnóstico más utilizados en el desarrollo de software paralelo.

\begin{table}[h!]
\centering
\renewcommand{\arraystretch}{1.5}
\caption{Parámetros opcionales comúnmente utilizados con \texttt{mpicc}.}
\label{tab:mpicc_flags}
\begin{tabular}{|l|p{11cm}|}
\hline
\textbf{Parámetro} & \textbf{Descripción y propósito} \\
\hline
\texttt{-g} & Genera información adicional de depuración en el ejecutable, permitiendo el análisis del programa con herramientas como GDB \cite{pacheco_2011}. \\
\hline
\texttt{-Wall} & Activa la emisión de advertencias detalladas sobre posibles inconsistencias en el código fuente \cite{pacheco_2011}. \\
\hline
\texttt{-o <outfile>} & Especifica el nombre del archivo ejecutable resultante \cite{pacheco_2011}. \\
\hline
\texttt{-O2} & Solicita la optimización del código generado, siendo la opción recomendada para pruebas de rendimiento \cite{pacheco_2011}. \\
\hline
\texttt{-O0} & Desactiva toda optimización del compilador, útil para asegurar que el código no sea alterado durante la depuración \cite{pacheco_2011}. \\
\hline
\texttt{-showme} & (Open MPI) Muestra el comando completo de compilación con rutas de inclusión y enlace sin ejecutarlo \cite{open_mpi_2026}. \\
\hline
\texttt{-compile-info} & (MPICH) Exhibe exclusivamente las banderas y rutas de preprocesamiento del compilador base. \\
\hline
\texttt{-link-info} & (MPICH) Muestra las banderas de enlace y las bibliotecas dinámicas aplicadas al binario. \\
\hline
\end{tabular}
\end{table}


El uso de los parámetros descriptivos como \texttt{-showme}, \texttt{-compile-info} y \texttt{-link-info} resulta de crucial importancia cuando se requiere auditar las rutas exactas del sistema de archivos asignadas por la instalación de MPI \cite{open_mpi_2026}. Estas banderas revelan el comando real subyacente que se ejecutará, facilitando la extracción de las rutas necesarias para la configuración de \textit{Makefiles} u otras herramientas de construcción automatizada \cite{gropp_using_mpi_1999}. Asimismo, la inclusión de la bandera \texttt{-g} asegura que los metadatos de depuración no se pierdan durante la compilación, permitiendo el análisis del binario mediante depuradores paralelos comerciales o de código abierto.

El uso correcto de \texttt{mpicc} demanda una atención rigurosa respecto al aprovisionamiento del entorno del sistema y la coherencia de las variables globales de la terminal\footnote{Por \comillas{aprovisionamiento} se entiende preparar y configurar el entorno necesario antes de compilar: instalar o cargar la implementación de MPI (p. ej. Open MPI o MPICH), activar los módulos del clúster cuando proceda (p. ej. mediante \texttt{module load}), asegurarse de que las rutas binarias estén en \texttt{PATH} y que las librerías compartidas sean localizables mediante \texttt{LD\_LIBRARY\_PATH}/\texttt{LIBRARY\_PATH}; además de verificar que los directorios de inclusión contienen \texttt{mpi.h} y otras cabeceras requeridas \cite{pacheco_2011,open_mpi_2026}.}. Uno de los fallos más recurrentes al invocar este comando se manifiesta mediante el error de sistema operativo que indica que el comando no ha sido encontrado, lo cual generalmente denota una ausencia de carga de los módulos de entorno del clúster o una incorrecta adición de las rutas binarias en la variable de entorno \texttt{\$PATH} \cite{pacheco_2011}. La resolución de este problema técnico implica comúnmente el uso de gestores de entorno como \texttt{environment-modules} para asegurar que las dependencias requeridas se encuentren activas en la sesión de terminal actual.

Además de la correcta invocación sintáctica, la utilización de \texttt{mpicc} exige atención rigurosa frente a errores comunes del proceso de construcción. En primer lugar, dado que el \textit{wrapper} está diseñado para simplificar la invocación del compilador subyacente, intentar compilar un programa de MPI usando directamente el compilador base (como \texttt{gcc}) sin replicar manualmente los argumentos de inclusión y enlace resultará en errores de resolución de dependencias, como la incapacidad de encontrar \texttt{mpi.h} o referencias indefinidas a las funciones de MPI \cite{pacheco_2011, open_mpi_2026}. Por otro lado, es necesario evitar estrictamente el uso de opciones del compilador que alteren el tamaño o la disposición en memoria de los tipos de datos de C o que modifiquen las convenciones de llamada, ya que esto rompería la compatibilidad binaria (ABI) y podría generar fallos críticos en tiempo de ejecución \cite{mpi_forum_2025}.

Otra consideración crítica surge durante la fase de enlazado, pues los enlazadores en sistemas Unix operan tradicionalmente en una sola pasada \cite{gropp_using_mpi_1999}. Esto implica que el orden en el que se pasan las bibliotecas adicionales es vital; por ejemplo, si se integran bibliotecas de \textit{profiling} que interceptan funciones de MPI mediante la interfaz \texttt{PMPI}, la biblioteca del \textit{profiler} debe enlazarse antes que las bibliotecas estándar (por ejemplo, \texttt{-lmyprof -lpmpi -lmpi}) \cite{gropp_using_mpi_1999}. Si se listan en el orden incorrecto, el programa compilará sin errores aparentes, pero fallará silenciosamente en su tarea de perfilar la aplicación \cite{gropp_using_mpi_1999}.

Finalmente, resulta prioritario evitar la colisión semántica derivada de mezclar extensiones de código fuente o emplear implementaciones de bibliotecas divergentes durante la construcción del proyecto. Intentar compilar código escrito en C++ utilizando \texttt{mpicc} en lugar de su contraparte específica \texttt{mpicxx} provocará fallas críticas de compilación debido a la incapacidad del compilador base para resolver el tipado y las clases del estándar moderno \cite{open_mpi_2026}. De igual modo, la mezcla accidental de archivos objeto generados por implementaciones distintas del estándar, como Open MPI y MPICH, causará fallos de enlazado indeterminados, por lo que se recomienda realizar una limpieza absoluta de los binarios previos antes de cambiar de entorno de ejecución \cite{mpi_forum_2025}. En aplicaciones heterogéneas que combinan módulos compilados con \texttt{mpicc} para C y con \texttt{mpif90} para Fortran, es necesario configurar opciones especiales durante la compilación y el enlazado para garantizar que la inicialización de MPI levante correctamente los entornos de cada lenguaje; de lo contrario, la aplicación puede omitir cargar el soporte para uno de los lenguajes si la biblioteca de ejecución correspondiente no se enlazó adecuadamente \cite{mpi_forum_2025}.

\subsection{Explicación del comando \texttt{mpirun} y \texttt{mpiexec}}

El comando \texttt{mpirun} constituye la herramienta fundamental empleada en entornos de computación de alto rendimiento para iniciar y gestionar aplicaciones distribuidas basadas en MPI. A diferencia de los programas secuenciales convencionales, las aplicaciones paralelas requieren un gestor de procesos coordinado que inicialice múltiples instancias del mismo binario a través de una arquitectura de hardware específica \cite{gropp_using_mpi_1999}. Este comando actúa como una interfaz de lanzamiento universal que abstrae la complejidad de la comunicación inicial entre los nodos de cómputo y el sistema operativo local \cite{pacheco_2011}.

Aunque históricamente la comunidad utilizaba diversas variantes de \texttt{mpirun}, el estándar MPI introdujo formalmente \texttt{mpiexec} para proveer un mecanismo uniforme y portable que evite la confusión derivada de opciones no estandarizadas \cite{mpi_forum_2025}. En la práctica moderna, ambos comandos suelen coexistir como sinónimos funcionales en la mayoría de las implementaciones, aunque \texttt{mpiexec} se prefiere cuando se busca el máximo apego a la especificación oficial. El propósito primordial de este mecanismo de lanzamiento consiste en establecer el entorno de ejecución necesario para que los procesos puedan identificarse y comunicarse entre sí desde el momento de su inicialización, asignando a cada uno un identificador único conocido como rango dentro del comunicador global \texttt{MPI\_COMM\_WORLD} \cite{mpi_forum_2025, grama_intro_parallel_2003}.

La invocación básica requiere la especificación explícita de la cantidad de tareas y la ruta del binario previamente compilado. En su forma más elemental, la sintaxis consiste en definir el parámetro de conteo de procesos seguido inmediatamente por el nombre del ejecutable, como en el comando \texttt{mpirun -np 4 ./programa\_paralelo}, el cual instruye al gestor de la biblioteca a instanciar exactamente cuatro copias independientes del software en los núcleos disponibles \cite{pacheco_2011}. Con el fin de flexibilizar la ejecución en arquitecturas distribuidas complejas, existen múltiples parámetros opcionales orientados a controlar la distribución y el mapeo detallado de la carga de trabajo, los cuales se resumen en el Cuadro~\ref{tab:parametros_mpirun}.

\begin{table}[h]
\centering
\caption{Parámetros opcionales recurrentes en los comandos \texttt{mpirun} y \texttt{mpiexec}.}
\label{tab:parametros_mpirun}
\begin{tabular}{|l|p{10cm}|}
\hline
\textbf{Parámetro} & \textbf{Descripción y propósito} \\ \hline
\texttt{-np <N>} / \texttt{-n <N>} & Especifica la cantidad total de procesos o instancias que se crearán para la aplicación distribuida. \\ \hline
\texttt{-hostfile <archivo>} & Define la ruta hacia un archivo de texto que contiene la lista de nodos y las ranuras de cómputo disponibles. \\ \hline
\texttt{-host <nombres>} & Indica los nodos o máquinas de la red donde específicamente se deben ejecutar los procesos. \\ \hline
\texttt{-arch <arquitectura>} & Permite solicitar nodos que posean una arquitectura de hardware específica (por ejemplo, \texttt{x86\_64}). \\ \hline
\texttt{-wdir <directorio>} & Establece el directorio de trabajo para los procesos instanciados. \\ \hline
\texttt{-configfile <archivo>} & Permite proporcionar un archivo de texto con configuraciones detalladas cuando los argumentos de línea de comandos son demasiados o complejos. \\ \hline
\texttt{-map-by <estrategia>} & Modifica la estrategia de asignación de procesos, permitiendo agruparlos por nodo, ranura o socket. \\ \hline
\texttt{-x <var>} / \texttt{--env <var>} & Exporta variables de entorno específicas del sistema local hacia los entornos de los nodos remotos. \\ \hline
\end{tabular}
\end{table}

El parámetro estructural \texttt{-hostfile} adquiere una relevancia crítica cuando se transiciona de un entorno de pruebas mononodo a un clúster de computadoras interconectadas por red, ya que distribuye de manera balanceada o secuencial las instancias del binario entre las máquinas físicas especificadas \cite{grama_intro_parallel_2003}. Por otro lado, modificadores avanzados como \texttt{-map-by} y \texttt{-x} garantizan el control granular sobre la afinidad del procesador y la propagación de contextos de ejecución homogéneos en sistemas heterogéneos \cite{open_mpi_2026}. Adicionalmente, el comando \texttt{mpiexec} permite la ejecución de aplicaciones heterogéneas bajo el modelo MPMD (\textit{Multiple Program, Multiple Data}), lo cual significa que se pueden iniciar simultáneamente varios programas distintos, o el mismo programa con diferentes argumentos, separando las especificaciones con el símbolo de dos puntos (\texttt{:}) en la línea de comandos \cite{mpi_forum_2025}.

Un aspecto didáctico de alta relevancia para quienes inician en MPI es el comportamiento no determinista de la salida estándar (\texttt{stdout}). Es erróneo asumir que las sentencias de impresión aparecerán en la pantalla ordenadas según el rango de cada proceso, pues los múltiples procesos compiten de forma asíncrona por el acceso al dispositivo de salida, generando una salida desordenada o interrumpida mutuamente \cite{pacheco_2011}. Para evitar esta confusión, se debe estructurar el código de manera que un solo proceso (por ejemplo, el rango 0) recolecte la información de los demás y se encargue exclusivamente de imprimirla en pantalla.

La implementación exitosa de entornos paralelos mediante \texttt{mpirun} o \texttt{mpiexec} demanda una supervisión constante sobre la compatibilidad de las versiones de software y la asignación realista de recursos físicos. Un error recurrente consiste en intentar ejecutar un número de procesos significativamente mayor al número de núcleos disponibles, práctica conocida como sobreasignación o \textit{oversubscription}, que provoca una degradación severa en el rendimiento debido a la latencia inducida por el cambio de contexto excesivo \cite{gropp_using_mpi_1999}. Asimismo, resulta imperativo vigilar la correspondencia exacta entre la implementación utilizada para compilar el código y la herramienta empleada para su ejecución; si se genera un binario con \texttt{mpicc} de MPICH e intentar ejecutarlo con \texttt{mpirun} de Open MPI, se desencadenarán fallos críticos de enlace en tiempo de ejecución o comportamientos indeterminados \cite{open_mpi_2026}. Por lo tanto, la validación previa de los módulos activos y el apego a la sintaxis estándar de \texttt{mpiexec} constituyen las mejores prácticas para mitigar fallas en la computación distribuida y preservar la portabilidad de los \textit{scripts} de lanzamiento \cite{pacheco_2011, mpi_forum_2025}.

\subsection{Significado del parámetro \texttt{-np} y sus variantes}

El parámetro \texttt{-np}, cuyo acrónimo proviene de la expresión inglesa \textit{number of processes} (número de procesos), constituye el argumento más crítico e indispensable al momento de invocar el gestor de ejecución \texttt{mpirun} o \texttt{mpiexec}. Su propósito primordial consiste en definir de manera explícita el grado de paralelismo inicial, especificando la cantidad exacta de instancias idénticas del binario que el entorno de ejecución distribuido debe inicializar de forma simultánea \cite{gropp_using_mpi_1999}. Este valor numérico altera directamente la arquitectura lógica de la aplicación, ya que determina el tamaño absoluto o cardinalidad del comunicador global primario, denotado formalmente dentro del estándar como \texttt{MPI\_COMM\_WORLD}\footnote{En MPI, la cardinalidad de un comunicador equivale a su tamaño, es decir, al número total de procesos que lo integran. En el caso de \texttt{MPI\_COMM\_WORLD}, esa cardinalidad coincide con el valor proporcionado a \texttt{-np} o \texttt{-n}; en consecuencia, cada proceso obtiene un rango único en el intervalo $[0, n-1]$ y puede consultar dicho tamaño mediante la rutina \texttt{MPI\_Comm\_size}. Esta relación entre cantidad de procesos, rangos y comunicador global es una de las bases del modelo SPMD en MPI \cite{mpi_forum_2025,pacheco_2011}.} \cite{mpi_forum_2025}.

En el aspecto operativo, la utilización de este modificador exige una asignación posicional estricta, colocándose inmediatamente después del comando de lanzamiento y precediendo a la ruta del archivo ejecutable. Por ejemplo, al ejecutar la instrucción \texttt{mpirun -np 16 ./simulacion}, la biblioteca subyacente coordina la creación de dieciséis procesos independientes, asignándole a cada uno de ellos un identificador entero único conocido como rango (\textit{rank}), el cual oscila en el intervalo cerrado $[0, n-1]$, donde $n$ representa el argumento entero provisto \cite{pacheco_2011}. Dependiendo de la distribución específica de MPI configurada en el clúster, este parámetro cuenta con variantes sintácticas equivalentes aceptadas por la industria, tales como \texttt{-n} o \texttt{--c}, las cuales preservan intacta la misma semántica funcional de aprovisionamiento de tareas \cite{open_mpi_2026}. No obstante, para mantener la máxima portabilidad y el apego estricto al estándar, se recomienda priorizar el uso de \texttt{-n}, ya que la documentación oficial de MPI explicita esta forma sobre la variante histórica \texttt{-np} \cite{gropp_using_mpi_1999}.

La selección del valor entero que acompaña a \texttt{-np} o \texttt{-n} demanda una planificación rigurosa basada en la topología del hardware y en el modelo de descomposición de datos del algoritmo paralelo. Un error metodológico ordinario consiste en fijar de forma arbitraria un número de procesos que no guarde coherencia con los requerimientos geométricos de la aplicación, como ocurre en algoritmos de reducción en árbol o mallas bidimensionales que exigen estrictamente potencias de dos o cuadrados perfectos para evitar desbalances de carga de trabajo \cite{grama_intro_parallel_2003}. Asimismo, si el entero provisto supera la cantidad de ranuras de procesamiento físicas especificadas en el archivo de configuración de nodos (\textit{hostfile}) sin haber habilitado explícitamente banderas de tolerancia a la sobreasignación\footnote{Estas banderas permiten solicitar al gestor de ejecución que acepte más procesos que ranuras físicas disponibles en los nodos asignados. En Open MPI, la opción habitual es \texttt{--oversubscribe}, la cual autoriza la sobreasignación y evita que el lanzamiento falle de inmediato por falta de ranuras. Su uso es válido en escenarios controlados, pero puede degradar el rendimiento porque varios procesos comparten el mismo núcleo o compiten por tiempo de CPU \cite{open_mpi_2026}.}, el gestor de procesos abortará la inicialización de inmediato, lo que en la práctica evidencia que la solicitud excede la capacidad asignada para ese lanzamiento y fuerza al usuario a decidir entre ajustar \texttt{-np} o aceptar la sobreasignación bajo control explícito \cite{open_mpi_2026}.

Como se mencionó anteriormente, además de la planificación del número de procesos, el estándar MPI ofrece mecanismos adicionales que otorgan flexibilidad al momento del lanzamiento. Es importante notar que, en circunstancias ordinarias, los modificadores \texttt{-np} y \texttt{-n} operan bajo una semántica de \comillas{todo o nada}: el gestor de procesos intenta asignar exactamente el número de procesos solicitados y, si no es posible encontrar suficientes ranuras disponibles en el clúster (incluso después de habilitar \texttt{--oversubscribe}), la ejecución falla de inmediato con un error de recursos \cite{open_mpi_2026}. 

En aplicaciones heterogéneas del tipo MPMD (\textit{Multiple Program, Multiple Data}), el parámetro puede utilizarse múltiples veces en una misma ejecución separando los bloques con el símbolo de dos puntos (\texttt{:}), lo que permite indicar una cantidad de procesos específica para cada programa; por ejemplo, \texttt{mpiexec -n 5 programa\_A : -n 10 programa\_B} \cite{mpi_forum_2025}. Sin embargo, cada bloque sigue siendo una solicitud rígida: si las ranuras no están disponibles, el lanzamiento aborta. De manera complementaria, cuando existe incertidumbre acerca de la disponibilidad de recursos en tiempo de ejecución o cuando se desea que el gestor aproveche dinámicamente los nodos disponibles sin fallar si no se alcanza el número exacto solicitado, el modificador \texttt{-soft} ofrece una semántica alternativa y más permisiva\footnote{La sintaxis del modificador \texttt{-soft} varía según la implementación. En Open MPI, se utiliza como \texttt{-npernode <min>:<max>} o directamente \texttt{-soft <rango>}, permitiendo especificar un rango mínimo y máximo de procesos permitidos. El gestor iniciará tantos procesos como pueda dentro de ese rango sin abortar si no se alcanza el máximo \cite{open_mpi_2026}.}: permite proporcionar al sistema un conjunto o rango de procesos permitidos, indicando a MPI que inicie el mayor número de procesos posible basándose en la disponibilidad real del clúster sin generar un error fatal \cite{mpi_forum_2025}. Esta estrategia es particularmente útil en entornos compartidos o con colas de trabajos dinámicas donde las ranuras pueden estar ocupadas parcialmente.

Un error lógico muy común entre los desarrolladores principiantes ocurre al interactuar con la salida estándar (\texttt{stdout}) inmediatamente después de incrementar el valor de \texttt{-np}. Es erróneo asumir que las sentencias de impresión aparecerán en la pantalla ordenadas según el rango de cada proceso, pues todos los procesos iniciados compiten simultáneamente por la consola, generando una salida impredecible o intercalada \cite{pacheco_2011}. Para evitarlo, la impresión siempre debe ser gestionada explícitamente desde el código, delegándola a un único proceso (comúnmente el rango 0). En ese sentido, la organización correcta del número de procesos no solo evita errores de lanzamiento, sino que también ayuda a mantener un comportamiento estable y predecible durante la ejecución paralela.

\subsection{Flujo de compilación y ejecución}

El ciclo de vida del desarrollo de software en entornos de paso de mensajes comprende un flujo estructurado que difiere significativamente del paradigma secuencial convencional. Esta secuencia operativa requiere una coordinación precisa entre la fase de traducción sintáctica y la fase de despliegue operacional dentro de la arquitectura de cómputo paralelo \cite{pacheco_2011}. Dado que MPI es una especificación de biblioteca y no un compilador en sí, este flujo se apoya en programas y comandos que envuelven el proceso estándar de desarrollo para simplificar el manejo de dependencias y el inicio del trabajo en un entorno distribuido \cite{mpi_forum_2025, open_mpi_2026}.

El propósito primordial de este flujo consiste en transformar las abstracciones lógicas del código fuente en C en un binario ejecutable altamente optimizado que pueda ser distribuido, instanciado y sincronizado eficientemente a través de múltiples nodos interconectados \cite{gropp_using_mpi_1999}. Esta progresión asegura que los canales de comunicación interna y las estructuras de datos globales del estándar queden correctamente cimentadas antes de iniciar cualquier cómputo a gran escala \cite{mpi_forum_2025}. De este modo, el flujo actúa como un puente entre la programación declarativa del desarrollador y la ejecución física en el clúster. Para lograrlo, la progresión desde las líneas de código fuente hasta la ejecución formal en paralelo se realiza mediante etapas secuenciales bien definidas que involucran a los scripts envolvedores y gestores de procesos de la biblioteca instalada \cite{grama_intro_parallel_2003}. Con la finalidad de ilustrar detalladamente la transición de los datos y el rol de cada componente de software, en la Tabla~\ref{tab:fases_flujo} se estructuran cronológicamente las fases requeridas dentro de este ecosistema tecnológico.

\begin{table}[h]
\centering
\caption{Etapas constitutivas del flujo de compilación y ejecución en MPI. Elaboración propia con base en \cite{pacheco_2011, mpi_forum_2025, grama_intro_parallel_2003}.}
\label{tab:fases_flujo}
\begin{tabular}{|c|l|l|p{6cm}|}
\hline
\textbf{Fase} & \textbf{Etapa} & \textbf{Herramienta} & \textbf{Transformación y operación} \\ \hline
1 & Preprocesamiento & \texttt{mpicc} & Se resuelven las directivas de inclusión como \texttt{\#include <mpi.h>}. \\ \hline
2 & Compilación & Compilador base & El compilador nativo (ej. GCC) traduce el código a lenguaje ensamblador y objetos. \\ \hline
3 & Enlazado & \texttt{ld} / \texttt{mpicc} & Se vinculan las bibliotecas estáticas o dinámicas del estándar MPI al binario final. \\ \hline
4 & Distribución & \texttt{mpirun} / \texttt{mpiexec} & El gestor propaga el contexto, lee el \textit{hostfile} y asigna los rangos lógicos a los nodos. \\ \hline
5 & Inicialización & \texttt{MPI\_Init} & Configura el entorno de MPI para cada proceso y conforma el comunicador global \texttt{MPI\_COMM\_WORLD}. \\ \hline
6 & Ejecución & Entorno MPI & Se instancian las tareas en los nodos y se realiza la computación paralela con comunicación entre procesos. \\ \hline
7 & Finalización & \texttt{MPI\_Finalize} & Libera los recursos de comunicación asignados y concluye formalmente el entorno paralelo. \\ \hline
\end{tabular}
\end{table}

Una vez que el archivo binario ejecutable ha sido generado con éxito por el script envolvedor, la responsabilidad total del flujo se transfiere al subsistema de lanzamiento en tiempo de ejecución \cite{gropp_using_mpi_1999}. En esta etapa, el comando \texttt{mpirun} o \texttt{mpiexec} —que actúan como gestores de procesos— toman el control operacional de la distribución y coordinación del binario en el clúster. Este componente interactúa dinámicamente con las capas de abstracción de red \cite{open_mpi_2026} estableciendo canales de comunicación entre los nodos asignados y, de ser el caso, con los planificadores de recursos del clúster (como Slurm o PBS) para negociar la disponibilidad de procesadores. Una vez obtenidos los recursos, el gestor mapea los rangos operacionales—es decir, asigna un identificador único (\textit{rank}) a cada copia del programa—de forma homogénea entre los procesadores asignados \cite{open_mpi_2026}. Este mapeo garantiza que cada proceso sepa quién es dentro de la estructura jerárquica de la ejecución paralela. Es en este punto donde el binario estático—un archivo que no cambia tras su compilación—cobra una naturaleza distribuida y coordinada, transformándose en múltiples instancias que se comunican entre sí de manera sincronizada.

Para que este flujo funcione sin problemas, las herramientas de compilación y lanzamiento deben pertenecer a la misma implementación de MPI; por ejemplo, utilizar \texttt{mpicc} y \texttt{mpiexec} provistos específicamente por Open MPI o por MPICH \cite{open_mpi_2026}. Compilar con los envoltorios de una implementación e intentar ejecutar el binario con el lanzador de otra generalmente resulta en errores de tiempo de ejecución y fallas críticas, dado que las bibliotecas de enlace y los formatos internos de los ejecutables difieren entre proyectos \cite{open_mpi_2026}. Asimismo, resulta fundamental asegurar la persistencia y sincronización de los archivos ejecutables dentro de sistemas de almacenamiento compartido distribuidos, tales como Network File System (NFS) o Lustre \cite{grama_intro_parallel_2003}. Si los nodos remotos asignados por el gestor de procesos no tienen acceso a la versión idéntica y actualizada del binario debido a retardos en la replicación o problemas de montaje en el sistema de archivos, el comando de ejecución distribuida arrojará errores de archivo no encontrado o, en el peor de los casos, ejecutará código obsoleto, comprometiendo la integridad del experimento numérico \cite{pacheco_2011}.

La culminación exitosa de este flujo integrado demanda la observancia rigurosa de consideraciones relativas a la homogeneidad de la arquitectura de hardware subyacente. En un clúster típico existen dos tipos de máquinas: los \textit{nodos de acceso frontal} (máquinas desde las cuales los usuarios se conectan e inician compilaciones) y los \textit{nodos de cálculo} (máquinas dedicadas exclusivamente a ejecutar tareas paralelas). Aunque frecuentemente comparten la misma arquitectura de procesador nominal, sus generaciones de procesador pueden diferir significativamente \cite{gropp_using_mpi_1999}.

Un riesgo crítico latente ocurre cuando se compila la aplicación en un nodo frontal utilizando optimizaciones agresivas específicas de la arquitectura mediante el parámetro \texttt{-march=native}. Esta bandera de compilación instruye al compilador para que genere instrucciones de máquina optimizadas exclusivamente para el procesador específico del nodo frontal—aprovechando características modernas como SIMD extendido—asumiendo que el hardware objetivo es idéntico \cite{gropp_using_mpi_1999}. Sin embargo, cuando posteriormente se intenta ejecutar ese binario optimizado en nodos de cálculo que poseen microarquitecturas de procesador heredadas o divergentes (por ejemplo, procesadores más antiguos sin soporte para esas instrucciones avanzadas), el sistema operativo no puede interpretar las instrucciones compiladas, desencadenando fallos catastróficos de instrucción ilegal—señal \texttt{SIGILL}—en tiempo de ejecución, lo que aborta inmediatamente todos los procesos paralelos \cite{gropp_using_mpi_1999}.

Si bien el estándar de MPI introdujo explícitamente el uso de \texttt{mpiexec} para estandarizar y hacer portable el mecanismo de arranque, la especificación intencionalmente no exige cómo deben los sistemas operativos asignar la memoria, gestionar la afinidad de procesos o coordinar internamente las tareas en el hardware subyacente. Este diseño deliberado deja tales detalles como responsabilidad exclusiva del entorno de ejecución particular (el sistema operativo, el planificador, la arquitectura del hardware) \cite{mpi_forum_2025, gropp_using_mpi_1999}. Mantener el apego riguroso a este flujo—desde la compilación sin optimizaciones arquitectónicas específicas, pasando por el uso consistente de herramientas de la misma implementación, hasta la invocación correcta de las funciones de inicialización (\texttt{MPI\_Init}) y finalización (\texttt{MPI\_Finalize})—garantiza que la aplicación sea altamente portable y ejecutable en una gran variedad de plataformas sin necesidad de recompilación \cite{mpi_forum_2025}.

\subsection{Ejemplo ilustrativo de una aplicación paralela en C}

La asimilación práctica de la biblioteca MPI se facilita significativamente mediante el análisis de un programa clásico de transmisión de mensajes, comúnmente denominado programa de saludos o \textit{greetings}. Este diseño de software elemental no solo permite verificar que la infraestructura de comunicación del clúster esté operando de manera correcta, sino que también introduce los fundamentos de la filosofía de paso de mensajes en sistemas distribuidos \cite{pacheco_2011}. A través de este modelo, el programador puede observar cómo un único código fuente es capaz de bifurcar su comportamiento lógico en tiempo de ejecución en función del identificador asignado a cada proceso independiente \cite{gropp_using_mpi_1999}.

El núcleo de este ejemplo práctico radica en la implementación del paradigma SPMD (\textit{Single Program, Multiple Data}), donde múltiples instancias ejecutan las mismas instrucciones sobre conjuntos de datos locales. Aunque MPI cuenta con más de un centenar de funciones, este código utiliza únicamente el subconjunto mínimo de seis rutinas fundamentales (\texttt{MPI\_Init}, \texttt{MPI\_Finalize}, \texttt{MPI\_Comm\_size}, \texttt{MPI\_Comm\_rank}, \texttt{MPI\_Send} y \texttt{MPI\_Recv}), las cuales son suficientes para escribir programas paralelos completamente funcionales \cite{grama_intro_parallel_2003, gropp_using_mpi_1999}. Para comprender a fondo el rol de cada una de las subrutinas invocadas en el programa, en la Tabla~\ref{tab:funciones_ejemplo} se detalla la semántica de las funciones esenciales de MPI utilizadas en la construcción de este binario. Luego, en la Figura~\ref{fig:codigo_greetings} se detalla la sintaxis completa en lenguaje C para llevar a cabo esta demostración fundamental.

\begin{table}[h]
	\centering
	\caption{Funciones fundamentales de la API de MPI utilizadas en el ejemplo. Elaboración propia con base en \cite{grama_intro_parallel_2003, gropp_using_mpi_1999, pacheco_2011, mpi_forum_2025}.}
	\label{tab:funciones_ejemplo}
	\begin{tabular}{|l|p{6.0cm}|p{5.0cm}|}
		\hline
		\textbf{Función MPI}                      & \textbf{Parámetros que recibe}                                                  & \textbf{Propósito y descripción operativa}                                                                                     \\ \hline
		\multirow{2}{*}{\texttt{MPI\_Init}}       & \texttt{int *argc}: Puntero al número de argumentos.                            & \multirow{2}{5.0cm}{Inicializa el entorno de ejecución de MPI y establece los contextos de comunicación iniciales.}            \\ \cline{2-2}
		                                          & \texttt{char ***argv}: Puntero al vector de argumentos.                         &                                                                                                                                \\ \hline
		\multirow{2}{*}{\texttt{MPI\_Comm\_size}} & \texttt{MPI\_Comm comm}: Comunicador a evaluar (ej. \texttt{MPI\_COMM\_WORLD}). & \multirow{2}{5.0cm}{Extrae el número total de procesos activos asociados al comunicador global especificado.}                  \\ \cline{2-2}
		                                          & \texttt{int *size}: Puntero entero donde se almacena el total de procesos.      &                                                                                                                                \\ \hline
		\multirow{2}{*}{\texttt{MPI\_Comm\_rank}} & \texttt{MPI\_Comm comm}: Comunicador a evaluar.                                 & \multirow{2}{5.0cm}{Obtiene el identificador entero único (rango) asignado al proceso que invoca la rutina.}                   \\ \cline{2-2}
		                                          & \texttt{int *rank}: Puntero entero donde se almacena el identificador único.    &                                                                                                                                \\ \hline
		\multirow{6}{*}{\texttt{MPI\_Send}}       & \texttt{const void *buf}: Dirección de memoria de los datos.                    & \multirow{6}{5.0cm}{Envía un mensaje de forma síncrona bloqueante hacia un destino utilizando una etiqueta específica.}        \\ \cline{2-2}
		                                          & \texttt{int count}: Cantidad de elementos en el búfer.                          &                                                                                                                                \\ \cline{2-2}
		                                          & \texttt{MPI\_Datatype datatype}: Tipo de dato MPI.                              &                                                                                                                                \\ \cline{2-2}
		                                          & \texttt{int dest}: Rango del proceso destino.                                   &                                                                                                                                \\ \cline{2-2}
		                                          & \texttt{int tag}: Etiqueta entera de control.                                   &                                                                                                                                \\ \cline{2-2}
		                                          & \texttt{MPI\_Comm comm}: Comunicador utilizado.                                 &                                                                                                                                \\ \hline
		\multirow{7}{*}{\texttt{MPI\_Recv}}       & \texttt{void *buf}: Búfer de memoria para almacenar datos.                      & \multirow{7}{5.0cm}{Recibe un flujo de datos desde un origen determinado, bloqueando el proceso hasta completar la operación.} \\ \cline{2-2}
		                                          & \texttt{int count}: Capacidad máxima de elementos.                              &                                                                                                                                \\ \cline{2-2}
		                                          & \texttt{MPI\_Datatype datatype}: Tipo de dato esperado.                         &                                                                                                                                \\ \cline{2-2}
		                                          & \texttt{int source}: Rango del proceso origen.                                  &                                                                                                                                \\ \cline{2-2}
		                                          & \texttt{int tag}: Etiqueta esperada.                                            &                                                                                                                                \\ \cline{2-2}
		                                          & \texttt{MPI\_Comm comm}: Comunicador utilizado.                                 &                                                                                                                                \\ \cline{2-2}
		                                          & \texttt{MPI\_Status *status}: Estructura con metadatos.                         &                                                                                                                                \\ \hline
		\texttt{MPI\_Finalize}                    & Ninguno (\texttt{void}).                                                        & Libera los recursos asignados por la biblioteca MPI y concluye formalmente el entorno paralelo.                                \\ \hline
	\end{tabular}
\end{table}

\begin{figure}[ht]
	\centering
	\includegraphics[width=0.85\textwidth]{Joni/greetings.png}
	\caption{Programa de saludos distribuido en MPI C. Elaboración propia con base en \cite{pacheco_2011}.}
	\label{fig:codigo_greetings}
\end{figure}


La estructura expuesta en la Figura~\ref{fig:codigo_greetings} revela el flujo de control dual que caracteriza a las aplicaciones distribuidas controladas por rangos numéricos. Mientras que los procesos con un rango superior a cero empaquetan y transmiten una cadena de caracteres hacia el nodo central, el proceso con rango cero se configura en un estado de escucha activa para recopilar y desplegar secuencialmente la información en la consola de salida estándar \cite{pacheco_2011}. El bloque condicional (\texttt{if... else}) demuestra a la perfección el modelo de ejecución SPMD, donde un único código fuente genera comportamientos distintos basándose en el identificador o rango (\texttt{my\_rank}) de cada proceso \cite{pacheco_2011}. Además, expone de manera clara el concepto de memoria distribuida: la variable \texttt{greeting} existe de forma privada e independiente en la memoria de cada proceso, por lo que la única manera de compartir su contenido con el proceso maestro es a través de la red usando un envío explícito \cite{pacheco_2011}.


En la Figura~\ref{fig:salida_greetings} se presenta la captura la impresión de salida del programa, donde se detalla el comando de invocación y la impresión secuencial de los mensajes de saludo transmitidos por cada uno de los rangos operacionales. El desarrollo de la salida impresa evidencia la ejecución del comando de lanzamiento \texttt{mpirun} configurado con un entorno de cuatro procesos en paralelo a través del parámetro estructural \texttt{-np 4}. Como resultado de esta instrucción, el proceso identificado con el rango cero toma inmediatamente el control de la consola estándar imprimiendo su propio saludo de forma local. De manera subsiguiente, las líneas de texto correspondientes a los rangos uno, dos y tres aparecen de forma ordenada debido al bucle síncrono de recepción soplado en el código maestro, el cual restringe la lectura consecutiva de las memorias buffer mediante llamadas deterministas a la primitiva \texttt{MPI\_Recv} \cite{open_mpi_2026}.


\begin{figure}[H]
	\centering
	\includegraphics[width=0.8\textwidth]{Joni/salida_greetings.png}
	\caption{Captura de pantalla de la terminal con la salida del programa \textit{greetings} en MPI. Elaboración propia.}
	\label{fig:salida_greetings}
\end{figure}


La correcta ejecución de este ejemplo de saludos demuestra la capacidad de las primitivas de comunicación punto a punto para establecer orden y sincronía en sistemas de memoria distribuida. El acoplamiento exacto entre las llamadas a \texttt{MPI\_Send} ejecutadas por los trabajadores y el bucle iterativo de \texttt{MPI\_Recv} coordinado por el nodo maestro asegura que los mensajes transiten por el tejido de red sin pérdidas ni colisiones de datos \cite{grama_intro_parallel_2003}. Este patrón de comunicación asimétrica constituye la base conceptual sobre la cual se cimientan algoritmos de descomposición de dominios más complejos en la computación científica actual \cite{mpi_forum_2025}.

A pesar de la aparente simplicidad matemática de este programa, es imperativo tomar precauciones rigurosas para evitar condiciones de bloqueo mutuo o \textit{deadlocks} durante la fase de escalamiento en supercomputadoras. Si el desarrollador altera el orden de las llamadas lógicas o si el proceso receptor intenta leer datos desde un rango inexistente, la aplicación completa quedará congelada indefinidamente en la memoria del sistema, desperdiciando valiosos recursos de procesamiento \cite{gropp_using_mpi_1999}. Por lo tanto, el diseño meticuloso de la topología lógica de los datos y el uso correcto de las estructuras de control representan las mejores directrices para la ingeniería de software de alto rendimiento \cite{open_mpi_2026}. Este programa es, en definitiva, el ejemplo por excelencia porque ilustra el ciclo de vida completo de una aplicación de MPI, enmarcando las comunicaciones estrictamente entre los llamados de inicio y fin, y comprender este patrón básico de comunicación y ramificación es el paso indispensable antes de avanzar a operaciones colectivas o arquitecturas de topologías virtuales \cite{grama_intro_parallel_2003, gropp_using_mpi_1999}.
